\centerline{\bf \Huge Ориентированные графы}

\begin{flushright}
        \scriptsize{Эпизод 13. Нападение ангела}
\end{flushright}

\textbf{Определение.} Ориентированный граф называется \textit{сильно связным}, если для любой пары вершин $u$ и $v$ существует путь по ребрам графа из $u$ в $v$ и путь из $v$ в $u$. \textit{Компонентой сильной связности} ориентированного графа называется сильно связный подграф, не являющийся подграфом большего сильно связного подграфа.

\problem{1}
(a) Докажите, что вершины любого ориентированного графа единственным образом разбиваются на не пересекающиеся компоненты сильной связности.

\noindent
(b) В ориентированном графе с $n$ вершинами нет ориентированных циклов. Докажите, что его вершины можно пронумеровать числами $1,2, \ldots, n$ так, что любое ребро ведет от вершины с меньшим номером в вершину с большим номером.

\noindent
(c) Пусть вершины ориентированного графа разбиты на $k$ компонент сильной связности. Тогда их можно пронумеровать числами $1,2, \ldots, k$ так, что любое ребро исходного графа либо ведет из компоненты с меньшим номером в компоненту с большим номером, либо лежит внутри одной компоненты.

\problem{2}
Докажите, что в полном ориентированном графе можно поменять направление не более одного ребра так, чтобы граф стал сильно связным.

\problem{3}
В стране 101 город. Некоторые города соединены дорогами с односторонним движением, причём в каждый город входит 40 дорог и из каждого города выходит 40 дорог. Докажите, что из каждого города можно добраться до любого другого, проехав не более чем по трём дорогам.

\textbf{Определение.}
\textit{Гамильтонов путь} --- это простой путь в графе, который проходит через каждую его вершину ровно один раз.

\problem{4}
(a) Докажите, что в полном ориентированном графе есть гамильтонов путь;

\noindent
(b) Докажите, что в сильно связном полном ориентированном графе любая вершина принадлежит циклу длины 3;

\newpage

\hspace{3mm}

\noindent
(c) Докажите, что в сильно связном полном ориентированном графе есть гамильтонов цикл;

\noindent
(d) Докажите, что в сильно связном полном ориентированном графе есть вершина, при удалении которой граф останется сильно связным.

\problem{5}
Докажите, что если в ориентированном графе нет ориентированных циклов, а в любом ориентированном пути не более $n$ ребер, то вершины графа можно разбить на $n+1$ группу так, чтобы внутри каждой группы не было ребер.

\problem{6}
В городе на каждую площадь выходит не менее трёх улиц. На улицах введено одностороннее движение так, что можно проехать с любой площади на любую другую. Докажите, что можно закрыть одну улицу так, что проезд сохранится.

\problem{7}
Дан связный граф, который при удалении любого ребра не теряет связность (граф без мостов). Докажите, что на всех его ребрах можно расставить ориентацию так, чтобы полученный ориентированный граф был сильно связным.